\documentclass[11pt]{article}
\usepackage[T1]{fontenc}
\usepackage[utf8]{inputenc}
\usepackage{lmodern}
\usepackage{amsmath,amssymb,amsthm,mathtools}
\usepackage[margin=1in]{geometry}
\usepackage{microtype,booktabs,longtable,array,enumitem}
\usepackage{xcolor}
\usepackage{hyperref}
\usepackage{fancyhdr}
\hypersetup{colorlinks=true,linkcolor=blue!45!black,citecolor=blue!45!black,
 urlcolor=blue!45!black,pdftitle={Surreal Markov Generators at Every Valuation Scale},
 pdfauthor={Research manuscript prepared for Vladimir Reshetnikov}}
\pagestyle{fancy}
\fancyhf{}
\fancyhead[L]{\small Surreal Markov hierarchies}
\fancyhead[R]{\small\thepage}
\renewcommand{\headrulewidth}{0.3pt}
\setlength{\headheight}{14pt}
\setlength{\emergencystretch}{2em}
\setlength{\parskip}{3pt}
\numberwithin{equation}{section}
\newtheorem{theorem}{Theorem}[section]
\newtheorem{proposition}[theorem]{Proposition}
\newtheorem{lemma}[theorem]{Lemma}
\newtheorem{corollary}[theorem]{Corollary}
\theoremstyle{definition}
\newtheorem{definition}[theorem]{Definition}
\newtheorem{example}[theorem]{Example}
\theoremstyle{remark}
\newtheorem{remark}[theorem]{Remark}
\newcommand{\R}{\mathbb R}
\newcommand{\C}{\mathbb C}
\newcommand{\Q}{\mathbb Q}
\newcommand{\Z}{\mathbb Z}
\newcommand{\No}{\mathbf{No}}
\newcommand{\F}{\mathcal F}
\newcommand{\Ocal}{\mathcal O}
\newcommand{\mcal}{\mathfrak m}
\newcommand{\one}{\mathbf 1}
\newcommand{\res}{\operatorname{res}}
\newcommand{\lc}{\operatorname{lc}}
\newcommand{\rank}{\operatorname{rank}}
\newcommand{\tr}{\operatorname{tr}}
\newcommand{\supp}{\operatorname{supp}}
\newcommand{\diag}{\operatorname{diag}}
\newcommand{\im}{\operatorname{im}}
\newcommand{\adj}{\operatorname{adj}}
\newcommand{\Rea}{\operatorname{Re}}
\newcommand{\ord}{\operatorname{ord}}
\newcommand{\Span}{\operatorname{span}}
\newcommand{\Id}{\mathrm I}
\newcommand{\bigO}[1]{O_v(#1)}
\newcommand{\revsha}{\texttt{4cf691c7d951e037739d32d9f5c387dcce724f3c}}
\newcommand{\doi}[1]{\href{https://doi.org/#1}{\texttt{#1}}}

\title{\textbf{Surreal Markov Generators\\at Every Valuation Scale}\\[0.5em]
\large Stochastic retracts, prescribed effective chains,\\
relative stability, and rank-one realization}
\author{Research manuscript\\Prepared for Vladimir Reshetnikov}
\date{22 September 2026}
\begin{document}
\maketitle
\thispagestyle{empty}

\begin{abstract}
A finite matrix of positive surreal transition rates can encode several orders of
infinitesimal time scales, including scales that no single real parameter can
preserve globally. We study its normalized resolvent rather than assume a global
surreal matrix exponential. Directed rooted forests give every entry's valuation
and leading coefficient at every scale. The resulting real matrix shadows form
a finite nested hierarchy of stochastic projections; each transition between
successive projections is exactly the resolvent of an ordinary finite Markov
generator on a stochastic retract.

The main converse is a realization theorem: every compatible finite hierarchy of
stochastic projections \emph{and prescribed effective Markov generators} occurs
at any prescribed strictly ordered Hahn scales in an irreducible surreal rate
matrix. An intermediate-scale correction repairs positivity without introducing
an additional leading resolvent scale. We also prove an entrywise relative
stability bound, uniform over all positive resolvent parameters and without a
spectral-gap loss, and show that the entire finite leading-term diagram admits a
rank-one monomial realization. Finally, the effective characteristic polynomial
recovers the leading surcomplex eigenvalue amplitudes, including tied scales and
nonreal eigenvalues. All proofs use finite algebra in a set-sized Hahn workspace;
no countable convergence in the full surreal topology is asserted.
\end{abstract}

\paragraph{Status.}
This manuscript takes the new-theorems alternative, not the claimed solution of a
named published open problem. Classical forest identities, Markov-chain
aggregation, stochastic-idempotent structure, and rank-one metastability theory
are credited explicitly. The arbitrary-rank realization and stability package is
presented as a candidate original contribution after a targeted, non-exhaustive
comparison. Priority and absence from every repository archive are not certified.
The arguments are mathematical proofs, not Lean-checked proofs.

\clearpage
\begingroup
\small
\setlength{\parskip}{0pt}
\tableofcontents
\endgroup
\clearpage

\section{The question, the answer, and its boundaries}\label{sec:intro}

A real continuous-time Markov chain has a finite row Laplacian $L$: its
off-diagonal entries are nonpositive and $L\one=0$. Replacing its rates by positive
surreal numbers makes it possible to distinguish transitions of size
$\omega^{-1}$, $\omega^{-2}$, and $\omega^{-\omega}$ exactly. This raises a more
substantial question than whether familiar matrix identities still hold:
\begin{quote}
Which finite hierarchies of effective stochastic dynamics can occur when the
observation scale moves through an arbitrary ordered group of surreal magnitudes?
Can the hierarchy, its rare entries, and its surcomplex spectral amplitudes be
recovered without choosing a numerical small parameter?
\end{quote}
We answer this for the normalized resolvent
\begin{equation}\label{eq:intro-resolvent}
 R_L(s)=s(sI+L)^{-1},\qquad s>0.
\end{equation}
For an ordinary real chain, this is the transition matrix observed at an
independent exponential time of rate $s$. In a Hahn field we take
\eqref{eq:intro-resolvent} as a rational-algebraic definition. This choice avoids
assuming a globally defined surreal matrix exponential or a surreal-valued
probability measure on a path space.

Write $s=c\,t^\alpha$, where $\alpha$ belongs to the value group and $c>0$ is real.
The central object is
\[
 K_\alpha(c)=\res R_L(c\,t^\alpha),
\]
where residue means the coefficient of $t^0$ of a bounded Hahn series. Smaller
$s$, hence larger $\alpha$, corresponds to longer observation. A \emph{plateau}
is a range of $\alpha$ on which this matrix is constant in $c$. A
\emph{crossover} is a scale at which nontrivial dependence on $c$ remains.
These terms refer only to this precisely defined resolvent hierarchy.

\subsection{Main conclusions}

There are at most $n-1$ crossover scales for an irreducible $n$-state system.
Between them lie stochastic idempotents
\[
 I=P_0,\ P_1,\ldots,P_m=\one\pi_0,
 \qquad P_iP_j=P_jP_i=P_{\max(i,j)}.
\]
An idempotent need not be a blockwise averaging matrix: some states can distribute
mass among several surviving classes. The correct reduced state space is a
\emph{stochastic retract}, a factorization $P=UV$, $VU=I$ with both $U$ and $V$
row stochastic.

At each crossover the shadow is
\[
 K_\alpha(c)=U\,c(cI+G_\alpha)^{-1}V,
\]
for a real Markov row Laplacian $G_\alpha$. This matrix is determined by one
compressed shadow, $G_\alpha=(VK_\alpha(1)U)^{-1}-I$. The result makes no
reversibility, simple-eigenvalue, or unique-minimizing-forest assumption.

The converse is stronger than realizing an eigenvalue list. Prescribe the whole
nested family $P_j$ and, on each retract, any real effective generator whose
limiting projection is the next $P_j$. Theorem~\ref{thm:general-realization}
constructs an irreducible Hahn generator with exactly those crossovers at any
prescribed ordered scales. Thus the compatibility conditions in the forward
theorem are also sufficient.

Two further conclusions sharply delimit what this structure measures.
First, relative perturbations of all rates of valuation at least $\delta>0$
change \emph{every positive resolvent entry} by a relative error of valuation at
least $\delta$, uniformly in $s$. This includes entries whose ordinary residue is
zero. Second, every finite leading-term diagram can be reproduced by a real
one-parameter monomial model. Higher rank permits exact global distinctions, but
no finite diagram of this kind can by itself certify higher rank.

\subsection{What is classical, and what is being proposed as new}

The directed matrix-forest theorem is classical; our normalization is a change of
variable in the identities of Chebotarev--Agaev~\cite{CA}. Finite stochastic
idempotents also have a classical structure theory; a useful modern treatment is
Dyer--Greenhill--Ullrich~\cite{DGU}. These are inputs, not new discoveries.

Markov metastability is a substantial existing subject. In particular,
Gan--Cameron~\cite{GC} treat reversible and nonreversible chains, symmetric and
nonsymmetric cases, and graph-based critical time scales. Their discussion
places eigenvalue-exponent formulas in the Freidlin--Wentzell tradition.
Singular-perturbation limits and aggregation are also developed by
Yin--Zhang~\cite{YZ} and related work. Merely allowing nonreversible rates, or
recovering eigenvalue exponents from optimal forests in rank one, is not claimed
as new here. Their path-event and typical-transition-graph hierarchies can contain
more scales than the nonzero eigenvalue scales considered in this manuscript.
Our bound $n-1$ does not bound all those other notions of critical time.

Nested self-adjoint projections and prescribed hierarchical spectra have close
predecessors in hierarchical Laplacians~\cite{BK}. Generalized ultrametric
inverse $M$-matrices are another related matrix-theoretic family~\cite{NV}.
Consequently the reversible, partition-based special case of our inverse
construction is not proposed as a new general principle. The distinction here
is the unrestricted stochastic-retract setting, the simultaneous prescription
of each effective generator, and the positivity repair that does not change the
leading hierarchy.

Higher-rank Hahn specialization has an established tropical-geometric
context~\cite{JS}. We prove our finite sign-specialization lemma for completeness,
not as a new theorem about ordered vector spaces. Its application below preserves
all marked-forest leading terms and all effective matrices at once.

\begin{center}
\begin{tabular}{p{0.28\textwidth}p{0.64\textwidth}}
\toprule
Ingredient & Status in this manuscript\\
\midrule
Matrix-forest identities & Classical input, extended by polynomial identity.\\
Stochastic-idempotent splitting & Classical input, with a complete proof.\\
Forest spectral exponents & Classical rank-one ancestry; arbitrary-rank proof included.\\
Forward/reverse hierarchy theorem & Candidate contribution: exact compatible-crossover realization over arbitrary rank.\\
Uniform relative stability & Candidate contribution in its simultaneous, all-parameter and marked-entry formulation.\\
Finite rank-one realization & Candidate contribution for the full leading resolvent diagram; the separation lemma is classical.\\
\bottomrule
\end{tabular}
\end{center}

The literature comparison is targeted rather than exhaustive. It did not locate
the exact collection of realization, stability, and simultaneous specialization
statements proved here. That is a statement about the search performed, not a
proof of historical priority.

\subsection{Comparison with the Surreal repository}

The comparison was pinned to commit
\begin{center}\small\revsha\end{center}
The maintained catalogue contains extensive analytic, dynamical, and spectral
material~\cite{Repo}. In particular, its spectral report already treats finite
matrices, determinantal scales, perturbations, and regularization as a scale
filter. Those subjects are not reclaimed here. Our objects are row-stochastic
resolvent shadows, their generally nonorthogonal retracts, and the converse
realization of prescribed effective Markov dynamics.

The audit read the main documentation catalogue and the spectral report's
inventory, inventoried \texttt{docs/new}, and performed an indexed repository
search for \texttt{Markov}, which returned no matches. The 18 ZIP archives then in
\texttt{docs/new} were \emph{not} expanded and their full mathematical contents
were not searched. Thus the warranted claim is that the package was not found
in the inspected material, not an exhaustive assertion about every binary
archive. The accompanying \texttt{SOURCE\_AUDIT.md} records this boundary.

\section{Hahn workspaces and positive rate matrices}\label{sec:workspace}

\subsection{A set-sized setting inside the surreals}

Fix a nonzero, set-sized, divisible ordered abelian group $\Gamma$. Put
\[
 F=\R((t^\Gamma)),\qquad K=\C((t^\Gamma)).
\]
A Hahn series has well-ordered support in $\Gamma$, and multiplication is the
usual well-defined convolution. For $x\ne0$, set
\[
 v(x)=\min\supp x,\qquad \lc(x)=[t^{v(x)}]x,
 \qquad v(0)=+\infty.
\]
The order on $F$ is determined by the sign of its leading coefficient. Its
valuation ring and maximal ideal are
\[
 \Ocal=\{x:v(x)\ge0\},\qquad \mcal=\{x:v(x)>0\}.
\]
The residue map $\res:\Ocal\to\R$ takes the $t^0$ coefficient. For bounded
complex series its target is $\C$. A positive element of $F$ is bounded by a real
constant exactly when its valuation is nonnegative. All nonzero real numbers
have valuation zero.

The standard Hahn-field facts used below are that $F$ is real closed and $K$ is
algebraically closed when $\Gamma$ is divisible. The identification $K=F[i]$ is
coefficientwise. Conway normal form~\cite{Conway,Gonshor} supplies the connection
to surreal numbers: for an ordered subgroup $\Gamma\subseteq\No$, the map
\begin{equation}\label{eq:surreal-embedding}
 \sum_{\gamma}a_\gamma t^\gamma
 \longmapsto \sum_{\gamma}a_\gamma\omega^{-\gamma}
\end{equation}
is the corresponding normal-form embedding. The minus sign reverses the
increasing Hahn support into a decreasing Conway support.

Any finite collection of surreal or surcomplex coefficients is contained in one
such workspace: take the union of their normal-form supports, change signs, and
take the rational span. This is a set-sized divisible ordered subgroup of
$\No$. All subsequent operations take place in its Hahn field or, where stated,
a larger set-sized one. No set of all surreal numbers is formed.

\begin{remark}[No hidden convergence assertion]
Every determinant, forest sum, and matrix inverse in this paper is finite
algebra. Residue is a ring homomorphism, not a sequential limit in the fine
topology of $\No$. Limits in a real variable $c$ or $u$ later in the paper are
ordinary limits of finite real matrices. No assertion that $t^n$ tends to zero
in an arbitrary-rank Hahn topology is used.
\end{remark}

\subsection{Rates, leading data, and conventions}

Let $E$ be a strongly connected directed graph on $\{1,\ldots,n\}$, with $n\ge2$,
without loops. For every allowed edge $e=(i,j)$ choose $q_e\in F_{>0}$. Set
$q_{ij}=0$ on missing edges. The \emph{row Laplacian} is
\begin{equation}\label{eq:laplacian}
 L_{ij}=-q_{ij}\quad(i\ne j),\qquad
 L_{ii}=\sum_{j\ne i}q_{ij}.
\end{equation}
The Markov generator in the usual sign convention is $-L$. All matrices act on
column observables; distributions are row vectors. In particular $L\one=0$.

Each existing rate has a unique leading representation
\begin{equation}\label{eq:rates}
 q_e=a_e t^{\gamma_e}(1+\varepsilon_e),
 \qquad a_e>0\text{ real},\quad \gamma_e\in\Gamma,
 \quad v(\varepsilon_e)>0,
\end{equation}
allowing $\varepsilon_e=0$. The coefficients beyond $a_e$ need not be positive.
The leading data are the finite family $(a_e,\gamma_e)_{e\in E}$.

For $\rho\in\Gamma$ write $x=\bigO{\rho}$ to mean $v(x)\ge\rho$; for a matrix
this means the bound entrywise. This is a valuation assertion, not a bound with
an unspecified real constant. Products of finitely many factors
$1+\bigO{\rho}$, for $\rho>0$, again have the form $1+\bigO{\rho}$.

\section{The finite forest algebra}\label{sec:forests}

An \emph{in-forest} is a spanning directed forest in which every nonroot vertex
has exactly one outgoing edge and following outgoing edges leads to a root.
Thus its edges point toward the roots. Let $\F_k$ be the finite set of such
forests with $k$ edges, and put $d=n-1$. For $f\in\F_k$, let $r_f(i)$ be the root
reached from vertex $i$ and let
\[
 q(f)=\prod_{e\in f}q_e.
\]
The empty forest has weight one. Strong connectivity guarantees that
$\F_d\ne\varnothing$, and deleting edges gives $\F_k\ne\varnothing$ for every
$0\le k\le d$.

Define the scalar and matrix forest coefficients
\begin{equation}\label{eq:forest-coefficients}
 \sigma_k=\sum_{f\in\F_k}q(f),\qquad
 F_k(i,j)=\sum_{\substack{f\in\F_k\\r_f(i)=j}}q(f).
\end{equation}
Every row of $F_k$ sums to $\sigma_k$, and every forest has $n-k$ roots. Hence
\begin{equation}\label{eq:forest-rows-trace}
 F_k\one=\sigma_k\one,\qquad
 \tr F_k=(n-k)\sigma_k,\qquad F_0=I,\quad\sigma_0=1.
\end{equation}

\begin{proposition}[Classical matrix-forest identity]\label{prop:forest}
For the convention \eqref{eq:laplacian},
\begin{align}
 \det(sI+L)&=sD(s),&D(s)&=\sum_{k=0}^d\sigma_k s^{d-k},\label{eq:forest-det}\\
 \adj(sI+L)&=N(s),&N(s)&=\sum_{k=0}^d F_k s^{d-k}.\label{eq:forest-adj}
\end{align}
Consequently, for every positive $s\in F$,
\begin{equation}\label{eq:forest-resolvent}
 R_L(s)=\frac{N(s)}{D(s)}.
\end{equation}
It is an entrywise positive row-stochastic matrix. In particular all its entries
belong to $\Ocal$ and admit real residues.
\end{proposition}

\begin{proof}
The determinant and adjugate identities are the directed matrix-forest theorem
\cite{CA}. They are polynomial identities with integer coefficients in the edge
weights and $s$. Thus their validity for real weights implies their validity in
any characteristic-zero coefficient field; alternatively one may use the usual
cycle-cancellation proof with formal weights. There is no specialization or
continuity assumption on Hahn series in this step.

For $s>0$, all terms in $D(s)$ are positive, so $sI+L$ is invertible and cancellation
of $s$ gives \eqref{eq:forest-resolvent}. Equation~\eqref{eq:forest-rows-trace}
gives row sum one. For every $i,j$, a spanning tree rooted at $j$ contributes to
$N(s)_{ij}$, so that entry is strictly positive. It is at most one because the
other entries in its row are nonnegative. This proves boundedness.
\end{proof}

\begin{remark}
We use the same formulas for a reducible row Laplacian when needed. Some maximal
forest coefficients may then be zero, but the determinant has a positive
$s^n$ term and the resolvent is still nonnegative and row stochastic for $s>0$.
Irreducibility is needed for strict positivity of every entry, not for the
rational identity itself.
\end{remark}

\subsection{Leading forest profiles}

For a forest define its valuation cost and amplitude by
\begin{equation}\label{eq:forest-leading}
 w(f)=\sum_{e\in f}\gamma_e,
 \qquad A(f)=\prod_{e\in f}a_e.
\end{equation}
For each edge count define
\begin{align}
 h_k&=\min_{f\in\F_k}w(f),\label{eq:profile-h}\\
 b_k&=\sum_{\substack{f\in\F_k\\w(f)=h_k}}A(f),\label{eq:profile-b}\\
 B_k(i,j)&=\sum_{\substack{f\in\F_k\\w(f)=h_k\\r_f(i)=j}}A(f).
 \label{eq:profile-B}
\end{align}
There is no cancellation among leading terms because all amplitudes are positive.
Thus
\begin{equation}\label{eq:profile-properties}
 v(\sigma_k)=h_k,\quad
 \lc(\sigma_k)=b_k>0,\quad
 B_k\one=b_k\one,\quad
 \tr B_k=(n-k)b_k.
\end{equation}
Although $B_k=\res(t^{-h_k}F_k)$, its zero entries may conceal higher-valuation
marked forests. They will be recovered separately rather than discarded.

\section{Every scale and every marked entry}\label{sec:leading}

For $\alpha\in\Gamma$ define
\begin{align}
 m_\alpha&=\min_{0\le k\le d}\{h_k+(d-k)\alpha\},\label{eq:m-alpha}\\
 I_\alpha&=\{k:h_k+(d-k)\alpha=m_\alpha\},\label{eq:active}\\
 D_\alpha(c)&=\sum_{k\in I_\alpha}b_kc^{d-k},\qquad
 N_\alpha(c)=\sum_{k\in I_\alpha}B_kc^{d-k}.
 \label{eq:initial-polynomials}
\end{align}
The expression $m_\alpha$ is also the minimum of
$w(f)+(d-|f|)\alpha$ over all forests. Its definition uses only ordered-group
addition and comparison; drawing a Newton polygon in a real plane is neither
necessary nor, at arbitrary rank, a literal interpretation.

\begin{theorem}[The complete leading-forest formula]\label{thm:leading}
For every $\alpha\in\Gamma$ and real $c>0$,
\begin{equation}\label{eq:shadow-formula}
 K_\alpha(c):=\res R_L(c t^\alpha)
 =\frac{N_\alpha(c)}{D_\alpha(c)}.
\end{equation}
This is a real, nonnegative, row-stochastic rational matrix function of $c$.
More generally, every positive $s\in F$ has a unique form
$s=c t^\alpha(1+\eta)$ with $c>0$ real and $v(\eta)>0$, and its resolvent has the
same residue $K_\alpha(c)$.

For a fixed entry define
\begin{align}
 m_{ij,\alpha}
 &=\min_{\substack{f\in\bigcup_k\F_k\\r_f(i)=j}}
       \{w(f)+(d-|f|)\alpha\},\label{eq:marked-min}\\
 A_{ij,\alpha}(c)
 &=\sum_{\substack{r_f(i)=j\\w(f)+(d-|f|)\alpha=m_{ij,\alpha}}}
       A(f)c^{d-|f|}.\label{eq:marked-amplitude}
\end{align}
Then
\begin{align}
 v\bigl(R_L(c t^\alpha)_{ij}\bigr)
   &=m_{ij,\alpha}-m_\alpha,\label{eq:marked-valuation}\\
 \res\!\left(t^{-(m_{ij,\alpha}-m_\alpha)}R_L(c t^\alpha)_{ij}\right)
   &=\frac{A_{ij,\alpha}(c)}{D_\alpha(c)}>0.
   \label{eq:marked-leading-coeff}
\end{align}
In particular the data $(a_e,\gamma_e)$ determine the leading term of every
entry, including entries whose unnormalized residue is zero.
\end{theorem}

\begin{proof}
A denominator forest term $q(f)s^{d-|f|}$ has valuation
$w(f)+(d-|f|)\alpha$ and leading coefficient $A(f)c^{d-|f|}$. A finite sum of
positive Hahn elements has valuation the minimum of their valuations, and its
leading coefficient is the sum of the positive leading coefficients at that
minimum. Therefore
\[
 \res(t^{-m_\alpha}D(c t^\alpha))=D_\alpha(c)>0,
 \qquad
 \res(t^{-m_\alpha}N(c t^\alpha))=N_\alpha(c).
\]
The normalized denominator is a unit of $\Ocal$. Taking its inverse and then
residue proves \eqref{eq:shadow-formula}. The row-stochastic assertion follows
either from Proposition~\ref{prop:forest} or from
$N_\alpha(c)\one=D_\alpha(c)\one$.

The numerator for a marked entry is another nonempty finite positive forest
sum. Normalizing it at its own minimum gives
\eqref{eq:marked-valuation}--\eqref{eq:marked-leading-coeff}. Finally,
$(1+\eta)^{d-|f|}$ has residue one for every forest. Replacing $c t^\alpha$ by
$c t^\alpha(1+\eta)$ therefore changes none of the leading computations.
\end{proof}

\begin{proposition}[A valuation remainder certificate]\label{prop:remainder}
Let
\[
 \delta=\min_{e\in E}v(\varepsilon_e),
\]
where $\delta=+\infty$ if all corrections vanish. At a fixed $\alpha$, let
\[
 \Delta_\alpha=
 \min\{w(f)+(d-|f|)\alpha-m_\alpha:
                 w(f)+(d-|f|)\alpha>m_\alpha\},
\]
with minimum $+\infty$ for the empty set. Then, for every real $c>0$,
\begin{equation}\label{eq:remainder}
 R_L(c t^\alpha)-K_\alpha(c)
   =\bigO{\min(\delta,\Delta_\alpha)}.
\end{equation}
The valuation lower bound is independent of $c$. If both quantities are infinite,
the identity is exact.
\end{proposition}

\begin{proof}
Every forest weight is $A(f)t^{w(f)}(1+\bigO{\delta})$. After division by
$t^{m_\alpha}$, nonminimal leading forest terms have valuation at least
$\Delta_\alpha$ and corrections to minimal terms have valuation at least
$\delta$. Thus the normalized numerator and denominator equal their real initial
polynomials plus errors of valuation at least their minimum. The real denominator
$D_\alpha(c)$ is nonzero, so inversion of this unit preserves the error bound.
All powers of the nonzero real $c$ have valuation zero.
\end{proof}

This certificate distinguishes two sources of uncertainty: subleading terms
inside an edge rate, and nonminimal forests. Neither a metric on the whole
surreal class nor a convergence rate in an ordinary norm appears in it.

\section{The finite flag of stochastic projections}\label{sec:flag}

Call $\alpha$ \emph{critical} if $I_\alpha$ contains at least two indices. There
are only finitely many critical values: every possible one belongs to
\begin{equation}\label{eq:candidate-scales}
 \left\{\frac{h_\ell-h_k}{\ell-k}:0\le k<\ell\le d\right\}.
\end{equation}
Division by a positive integer is available because $\Gamma$ is divisible.
Away from these values the active index is unique and
$K_\alpha(c)=B_k/b_k$ is independent of $c$.

\begin{lemma}[Resolvent compatibility]\label{lem:resolvent}
For positive $s\ne u$,
\begin{equation}\label{eq:resolvent-identity}
 R_L(s)R_L(u)=\frac{uR_L(s)-sR_L(u)}{u-s}.
\end{equation}
Consequently
\begin{align}
 K_\alpha(c)K_\alpha(e)
   &=\frac{eK_\alpha(c)-cK_\alpha(e)}{e-c}
       &&(c,e>0,\ c\ne e),\label{eq:same-scale}\\
 K_\alpha(c)K_\beta(e)
   &=K_\beta(e)K_\alpha(c)=K_\beta(e)
       &&(\alpha<\beta).\label{eq:cross-scale}
\end{align}
\end{lemma}

\begin{proof}
The two resolvents commute since they are rational functions of $L$.
Subtracting $(sI+L)^{-1}$ and $(uI+L)^{-1}$ proves
\eqref{eq:resolvent-identity}. Substituting $s=ct^\alpha$, $u=et^\alpha$ and
taking residues gives \eqref{eq:same-scale}. If $\alpha<\beta$, divide numerator
and denominator on the right of \eqref{eq:resolvent-identity} by $s$. The ratio
$u/s$ has positive valuation; all resolvent entries are bounded. The residue of
$((u/s)R_L(s)-R_L(u))/((u/s)-1)$ is therefore $K_\beta(e)$.
\end{proof}

\begin{theorem}[Projection hierarchy]\label{thm:flag}
There are critical scales $\alpha_1<\cdots<\alpha_m$, with $1\le m\le n-1$,
and distinct stochastic idempotents
\begin{equation}\label{eq:flag}
 P_0=I,\ P_1,\ldots,P_m=\one\pi_0,
 \qquad P_iP_j=P_jP_i=P_{\max(i,j)},
\end{equation}
where $\pi_0$ is a real probability row vector, possibly with zero entries.
On the open interval between $\alpha_j$ and $\alpha_{j+1}$ the shadow is $P_j$;
the initial and final intervals have shadows $I$ and $\one\pi_0$.
The ranks decrease strictly.

At a critical scale $\alpha$, write
\[
 k_-=\min I_\alpha,\qquad k_+=\max I_\alpha,
 \qquad p=n-k_-,\qquad q=n-k_+.
\]
The preceding and following plateaux satisfy
\begin{align}
 P&=\lim_{c\to\infty}K_\alpha(c)=B_{k_-}/b_{k_-},\label{eq:fast-endpoint}\\
 Q&=\lim_{c\downarrow0}K_\alpha(c)=B_{k_+}/b_{k_+},\label{eq:slow-endpoint}
\end{align}
with $\rank P=p$, $\rank Q=q$, and
\begin{equation}\label{eq:endpoint-relations}
 PK_\alpha(c)=K_\alpha(c)P=K_\alpha(c),\qquad
 K_\alpha(c)Q=QK_\alpha(c)=Q.
\end{equation}
Both displayed limits are limits of real rational matrices.
\end{theorem}

\begin{proof}
At a noncritical scale $K_\alpha(c)=P$ for every $c$. Taking two different real
values in \eqref{eq:same-scale} gives $P^2=P$. If $k$ is the active index,
\eqref{eq:profile-properties} gives
\[
 \rank P=\tr P=n-k,
\]
since the eigenvalues of an idempotent are zero and one. Cross-scale
compatibility proves all products in \eqref{eq:flag}.

As $\alpha$ increases, a lower envelope of lines with slopes $d-k$ can only move
to larger $k$. More explicitly, if $k$ is active at $\alpha$ and $\ell$ is active
at $\beta>\alpha$, adding their two minimizing inequalities gives
$(\ell-k)(\beta-\alpha)\ge0$. Thus $\ell\ge k$.
Choosing $\alpha$ below all the finitely many crossing values makes $k=0$ the
unique minimizer; choosing it above all of them makes $k=d$ the unique minimizer.
These choices exist in any nonzero ordered group. The first matrix is $I$. The
last has rank one, is stochastic, and so has identical probability rows.

At a critical value, the highest power of $c$ in both initial polynomials comes
from $k_-$ and the lowest comes from $k_+$. This proves the endpoint limits.
Immediately before and after that value, the corresponding unique minimizers
are $k_-$ and $k_+$. Since $k_-<k_+$, the rank drops by a positive integer at
each critical value. The total rank loss is $n-1$, proving $m\le n-1$.
Finally, apply \eqref{eq:cross-scale} using any scale in the adjacent plateau
intervals to obtain \eqref{eq:endpoint-relations}.
\end{proof}

\begin{remark}[What the last row vector means]
The exact stationary distribution is given by the tree formula
$\pi_j=F_d(i,j)/\sigma_d$, independently of $i$. Hence $\pi_0=\res\pi$.
Exact irreducibility makes every $\pi_j$ positive, but does not force every
$\res\pi_j$ to be positive. Infinitesimally occupied states explain the zero
columns and transient mixing that can occur in intermediate projections.
\end{remark}

\section{Every crossover is a real effective Markov chain}\label{sec:effective}

\subsection{Stochastic retracts, including transient states}

\begin{lemma}[Classical splitting of a stochastic idempotent]\label{lem:splitting}
If $P$ is an $n\times n$ real row-stochastic idempotent of rank $p$, there exist
nonnegative row-stochastic matrices
\[
 U\in\R^{n\times p},\qquad V\in\R^{p\times n},
 \qquad P=UV,\qquad VU=I_p.
\]
A canonical choice, up to ordering, is indexed by the closed communicating
classes of $P$: the rows of $V$ are their stationary probability vectors, extended
by zero, and $U_{ia}$ is the mass that row $i$ assigns to class $a$.
\end{lemma}

\begin{proof}
Here is a direct proof of the standard structure result~\cite{DGU}. Decompose
the directed graph of positive entries of $P$ into its closed communicating
classes $C_a$ and its remaining states $T$. With closed classes first, its matrix
has the form
\[
 P=\begin{pmatrix}C&0\\ B&T_0\end{pmatrix},
\]
where $C$ is block diagonal with irreducible stochastic blocks. Every transient
state has a finite positive-probability path to a closed class. By finiteness,
there are an integer $r$ and a real $\eta>0$ such that every row sum of $T_0^r$
is at most $1-\eta$. It follows that $T_0^m\to0$ as $m\to\infty$.
But $P^2=P$ implies $T_0^2=T_0$, so $T_0=0$.

For a closed irreducible block $C_a$, each real vector $x$ with $C_ax=x$ is
constant: choose a maximal coordinate; its expression as a convex combination
forces the same maximum at every successor, and irreducibility propagates it to
all coordinates. Every column of $C_a$ is fixed by $C_a$, since $C_a^2=C_a$.
Therefore all its columns are constant, and
$C_a=\one\nu_a$ for a probability row vector $\nu_a$. Its entries are positive
by irreducibility. The equation $P^2=P$ now gives $B=BC$, so the portion of any
transient row on $C_a$ is a nonnegative scalar multiple of $\nu_a$.

Extend each $\nu_a$ by zero to a row of $V$, and set
$U_{ia}=\sum_{j\in C_a}P_{ij}$. The preceding description proves $UV=P$.
Both matrices have row sums one. For $i\in C_a$, row $i$ of $U$ is the $a$th
coordinate row, so $VU=I$. These identities imply that the number of classes is
$\rank P=p$.
\end{proof}

The rows of $U$ at transient states need not be coordinate rows. Replacing this
factorization by a partition quotient would discard real information and would
make the next theorem false in its intended generality.

\subsection{Reconstructing the effective generator}

\begin{theorem}[Effective-generator reconstruction]\label{thm:effective}
Fix a critical $\alpha$, with endpoints $P,Q$ of ranks $p>q$, and take the
canonical splitting $P=UV$ from Lemma~\ref{lem:splitting}. Set
\[
 H(c)=VK_\alpha(c)U.
\]
There is a unique real $p\times p$ row Laplacian $G_\alpha$ such that
\begin{equation}\label{eq:effective-resolvent}
 H(c)=c(cI_p+G_\alpha)^{-1},\qquad
 K_\alpha(c)=U\,c(cI_p+G_\alpha)^{-1}V
 \quad(c>0).
\end{equation}
It is given explicitly by
\begin{equation}\label{eq:effective-one}
 G_\alpha=H(1)^{-1}-I_p.
\end{equation}
Its zero eigenvalue is semisimple, its kernel has dimension $q$, and its
zero-eigenvalue projection is $\widehat Q=VQU$. In particular
$\rank G_\alpha=p-q$. Ordering the closed classes differently only permutes
$G_\alpha$.
\end{theorem}

\begin{proof}
Because $U,V,K_\alpha(c)$ are stochastic, so is $H(c)$. Equations
\eqref{eq:endpoint-relations} give $K_\alpha(c)=UH(c)V$, while
\eqref{eq:same-scale} gives
\begin{equation}\label{eq:compressed-identity}
 H(c)H(e)=\frac{eH(c)-cH(e)}{e-c}.
\end{equation}
Also $H(c)\to VPU=I_p$ as $c\to\infty$. Choose a sufficiently large real $e$
so that $H(e)$ is invertible. Rearranging \eqref{eq:compressed-identity} yields
\[
 H(c)\bigl[cI_p+e(H(e)^{-1}-I_p)\bigr]=cI_p
\]
for $c\ne e$, and the identity extends to $c=e$. Define
$G_\alpha=e(H(e)^{-1}-I_p)$. A square matrix with a right inverse is invertible,
so the first identity of \eqref{eq:effective-resolvent} holds for every $c>0$.
It then gives \eqref{eq:effective-one} and uniqueness.

Since $H(c)\one=\one$, its inverse also fixes $\one$, and $G_\alpha\one=0$.
For large real $c$,
\[
 H(c)=I_p-c^{-1}G_\alpha+O(c^{-2}).
\]
Every off-diagonal entry of $H(c)$ is nonnegative. Multiplication by $c$ and
passage to the real limit therefore gives $(G_\alpha)_{ab}\le0$ for $a\ne b$.
Thus it is a real Markov row Laplacian.

The limit as $c\downarrow0$ is $\widehat Q=VQU$. It is an idempotent, and
$U\widehat QV=Q$, so its rank is $q$. From
$G_\alpha H(c)=c(I-H(c))$ and boundedness of $H(c)$, we obtain
$G_\alpha\widehat Q=0$. Conversely, if $G_\alpha x=0$, then $H(c)x=x$ for
all $c$ and hence $\widehat Qx=x$. Thus the kernel is exactly
$\im\widehat Q$. The same argument on the other side gives
$\widehat QG_\alpha=0$.

There cannot be a nontrivial Jordan chain $G_\alpha y=x\ne0$,
$G_\alpha x=0$. Such a chain would give
$H(c)y=y-x/c$, which is unbounded as $c\downarrow0$, contrary to the bounded
matrix $H(c)$. A longer nontrivial Jordan chain contains one of length two.
This proves semisimplicity and all remaining claims.
\end{proof}

\begin{corollary}[A real semigroup on each retract]\label{cor:effective-semigroup}
For real $u\ge0$ define
\[
 S_\alpha(u)=Ue^{-uG_\alpha}V.
\]
These matrices are stochastic,
$S_\alpha(0)=P$, $S_\alpha(u)S_\alpha(v)=S_\alpha(u+v)$, and
$S_\alpha(u)\to Q$ as $u\to\infty$. Their ordinary real Laplace transform is
\[
 \int_0^\infty c e^{-cu}S_\alpha(u)\,du=K_\alpha(c).
\]
\end{corollary}

\begin{proof}
Choose a real $M>\max_a(G_\alpha)_{aa}$. Then
$T=I-G_\alpha/M$ is stochastic, and the ordinary absolutely convergent identity
\[
 e^{-uG_\alpha}=e^{-Mu}\sum_{r\ge0}\frac{(Mu)^r}{r!}T^r
\]
proves nonnegativity and row sums one. The semigroup identity follows from
$VU=I$. To see the large-$u$ limit, all nonzero eigenvalues of a real row
Laplacian have strictly positive real part: if $Gz=\lambda z$ and $|z_a|$ is
maximal and nonzero, the $a$th row shows $\Rea\lambda\ge0$. Equality forces
$z_b/z_a=1$ for every positive outgoing rate, and then forces $\lambda=0$.
The zero eigenvalue is semisimple by the preceding theorem, so the exponential
converges to $\widehat Q$. The Laplace transform identity follows by integrating
the finite-dimensional real resolvent, or directly by differentiating
$e^{-cu}e^{-uG_\alpha}$ and evaluating its endpoints.
\end{proof}

\begin{remark}
This is a theorem about a \emph{real} effective semigroup. It is not a claim that
$S_\alpha(u)$ is the residue of a previously defined global surreal exponential
$\exp(-u t^{-\alpha}L)$. No such global analytic construction is required.
\end{remark}

\section{The converse: realizing a prescribed hierarchy}\label{sec:inverse}

This section is the principal inverse result. It first realizes arbitrary
projection flags explicitly, then realizes prescribed effective generators at
every transition. The extra difficulty in the second construction is
entrywise positivity, not spectral decomposition.

\subsection{Every stochastic projection flag occurs}

\begin{theorem}[Projection-flag realization]\label{thm:flag-realization}
Let $P_0=I,P_1,\ldots,P_m$ be distinct real stochastic idempotents satisfying
\[
 P_iP_j=P_jP_i=P_{\max(i,j)},\qquad \rank P_m=1.
\]
Choose any $\alpha_1<\cdots<\alpha_m$ in $\Gamma$, and write
$\tau_j=t^{\alpha_j}$. Then
\begin{equation}\label{eq:simple-inverse}
 H=\sum_{j=1}^m\tau_j(P_{j-1}-P_j)
\end{equation}
is a Hahn row Laplacian. Its normalized resolvent is exactly
\begin{equation}\label{eq:simple-inverse-resolvent}
 R_H(s)=P_m+\sum_{j=1}^m\frac{s}{s+\tau_j}(P_{j-1}-P_j).
\end{equation}
Its plateaux are precisely the prescribed $P_j$, and its crossover at $\alpha_j$
is
\begin{equation}\label{eq:simple-inverse-shadow}
 K_{\alpha_j}(c)=P_j+\frac{c}{c+1}(P_{j-1}-P_j).
\end{equation}
An irreducible row Laplacian with the same shadows at \emph{every} scale is
obtained by the completion in Lemma~\ref{lem:completion} below.
\end{theorem}

\begin{proof}
Put $E_j=P_{j-1}-P_j$. These are algebraic idempotents with
$E_iE_j=0$ for $i\ne j$, $E_jP_m=P_mE_j=0$, and
$P_m+\sum_jE_j=I$. Each $E_j$ has positive rank. They need not be nonnegative.
The alternative expression
\begin{equation}\label{eq:inverse-signs}
 H=\tau_1I+\sum_{j=1}^{m-1}(\tau_{j+1}-\tau_j)P_j-\tau_mP_m
\end{equation}
shows that every off-diagonal entry is nonpositive, because
$0<\tau_{j+1}<\tau_j$. Also $H\one=0$. Multiplying
\eqref{eq:simple-inverse-resolvent} by $sI+H$ proves the exact identity.

At scale $s=ct^{\alpha_j}$, the residue of $s/(s+\tau_i)$ is zero for $i<j$,
$c/(c+1)$ for $i=j$, and one for $i>j$. Telescoping the last terms gives
\eqref{eq:simple-inverse-shadow}. At a scale strictly between successive
$\alpha_j$, the same calculation gives $P_j$. Since $E_j\ne0$, no stated
crossover disappears and there are no others.
\end{proof}

\begin{lemma}[Irreducible completion invisible to the hierarchy]\label{lem:completion}
Suppose a Hahn row Laplacian $H$ has final shadow $P_m=\one\pi_0$ at all scales
strictly above $\alpha_m$. Choose $\beta>\alpha_m$, let $\epsilon=t^\beta$, and
let $J=\one\nu$ for any strictly positive real probability row $\nu$. Then
\begin{equation}\label{eq:completion}
 L=H+\epsilon(I-J)
\end{equation}
has strictly negative off-diagonal entries. If $H$ has the hierarchy of
Theorem~\ref{thm:flag-realization}, or of
Theorem~\ref{thm:general-realization} below, $L$ has the same shadow at every
$\alpha$ and every $c>0$.
\end{lemma}

\begin{proof}
Strict negativity follows from the added rate $\epsilon\nu_j>0$ on every edge.
For every positive $s$, a direct rank-one inverse computation gives
\begin{align}
 R_L(s)
  &=(sI+\epsilon J)((s+\epsilon)I+H)^{-1}\label{eq:completion-exact}\\
  &=\left(\frac{s}{s+\epsilon}I+
           \frac{\epsilon}{s+\epsilon}J\right)R_H(s+\epsilon).
           \label{eq:completion-convex}
\end{align}
For completeness, if $M=(s+\epsilon)I+H$, then $M^{-1}J=J/(s+\epsilon)$ because
$H\one=0$. Thus
$(M-\epsilon J)^{-1}=(I+\epsilon J/s)M^{-1}$, proving
\eqref{eq:completion-exact}. The order of these factors matters; no commutation
of $J$ and $H$ is assumed.

If $s=ct^\alpha$ and $\alpha<\beta$, the first factor in
\eqref{eq:completion-convex} has residue $I$, and $s+\epsilon$ has the same
leading term as $s$. The leading-forest argument, valid also for reducible
$H$, shows that the shadow is unchanged. If $\alpha\ge\beta$, the parameter
$s+\epsilon$ has valuation $\beta$, so its $H$-shadow is $P_m$.
The first factor has residue a stochastic combination of $I$ and $J$.
Since $JP_m=P_m$, the product still has residue $P_m$. In particular the new
edge scale $\beta$ creates no new leading crossover.
\end{proof}

\subsection{Admissible effective data}

Let a flag $P_0,\ldots,P_m$ as above be given. For each $j$, choose the canonical
splitting
\[
 P_{j-1}=U_jV_j,\qquad V_jU_j=I_{p_{j-1}},\qquad
 p_j=\rank P_j,
\]
and set
\begin{equation}\label{eq:macro-next}
 Q_j=V_jP_jU_j.
\end{equation}
It is a stochastic idempotent of rank $p_j$, and $U_jQ_jV_j=P_j$.

\begin{definition}[Compatible effective generators]\label{def:compatible}
A family of real row Laplacians $G_j$ of sizes $p_{j-1}$ is compatible with the
flag if
\begin{equation}\label{eq:compatible}
 \lim_{c\downarrow0}c(cI+G_j)^{-1}=Q_j.
\end{equation}
Equivalently, $Q_j$ is the zero-eigenvalue projection of $G_j$.
\end{definition}

There is always at least one compatible choice, namely $G_j=I-Q_j$. But the
next theorem allows any compatible $G_j$, including nonreversible generators,
complex eigenvalues, unequal eigenvalues within a single scale, and nontrivial
Jordan blocks at nonzero eigenvalues.

\begin{theorem}[Prescribed-crossover realization]\label{thm:general-realization}
Fix any flag $P_0=I,P_1,\ldots,P_m$ of distinct stochastic idempotents ending
in rank one, any compatible real generators $G_j$, and any prescribed scales
$\alpha_1<\cdots<\alpha_m$ in a nonzero divisible ordered group $\Gamma$.
There exists an irreducible Hahn row Laplacian $L$ whose only leading resolvent
crossovers are those scales, whose plateaux are the $P_j$, and whose shadows at
them are exactly
\begin{equation}\label{eq:prescribed-shadow}
 K_{\alpha_j}(c)=U_j\,c(cI_{p_{j-1}}+G_j)^{-1}V_j
 \qquad(c>0).
\end{equation}
It can be chosen with finite-support rates. Moreover, after replacing every
existing rate by its leading monomial, it can be chosen with positive monomial
rates while preserving all these shadows.
\end{theorem}

\begin{proof}
\emph{Step 1: lift the effective generators.}
Set
\[
 E_j=P_{j-1}-P_j,\qquad A_j=U_jG_jV_j.
\]
Compatibility implies $G_jQ_j=Q_jG_j=0$, and hence
\begin{equation}\label{eq:lifted-blocks}
 A_j=E_jA_jE_j,\qquad A_j\one=0.
\end{equation}
On the real subspace $\im E_j$, the operator $A_j$ is invertible. Indeed the
restriction is identified through $U_j,V_j$ with the restriction of $G_j$ to
$\ker Q_j$, where zero is not an eigenvalue. The mutually annihilating
idempotents $E_1,\ldots,E_m,P_m$ therefore decompose the whole vector space
into the intended spectral blocks. In particular $A_iA_j=0$ for $i\ne j$.

The naive sum $\sum_jt^{\alpha_j}A_j$ need not be a row Laplacian. A lifted
$A_j$ may have a positive off-diagonal entry on a pair of original states in the
same preceding class. The next step repairs this without altering a prescribed
crossover.

\emph{Step 2: insert positivity corrections between scales.}
For $2\le j\le m$ choose
$\alpha_{j-1}<\beta_j<\alpha_j$, possible by divisibility. Put
\begin{equation}\label{eq:general-construction}
 H=\sum_{j=1}^m t^{\alpha_j}A_j
       +\sum_{j=2}^m t^{\beta_j}(I-P_{j-1}).
\end{equation}
Every row sum is zero. We check the off-diagonal signs carefully.
If $(P_{j-1})_{ab}=0$, nonnegativity of $U_j,V_j$ implies
$(U_j)_{ar}(V_j)_{rb}=0$ for each $r$. Thus all diagonal contributions to
$(U_jG_jV_j)_{ab}$ vanish, and its remaining contributions are nonpositive.
Consequently
\begin{equation}\label{eq:sign-cover}
 (A_j)_{ab}>0\quad\Longrightarrow\quad (P_{j-1})_{ab}>0
 \qquad(a\ne b).
\end{equation}
For $j=1$ the preceding projection is $I$, so $A_1=G_1$ has no positive
off-diagonal entries. For $j\ge2$, any positive contribution
$t^{\alpha_j}(A_j)_{ab}$ is preceded by the strictly negative contribution
$-t^{\beta_j}(P_{j-1})_{ab}$ of smaller valuation. All exponents in the ordered
list
\[
 \alpha_1<\beta_2<\alpha_2<\cdots<\beta_m<\alpha_m
\]
are distinct. Therefore the earliest nonzero coefficient in any off-diagonal
entry of $H$ cannot be positive: every possible positive coefficient has an
earlier negative one. The leading coefficient is negative, or the entry is
zero. This proves that $H$ is a row Laplacian. It is a finite sum of monomials.

\emph{Step 3: show that the inserted scales are invisible.}
On $\im E_i$, the operator $I-P_{j-1}$ is the identity if $i<j$ and zero if
$i\ge j$. Thus the block of $H$ on $\im E_i$ is
\begin{equation}\label{eq:realization-block}
 t^{\alpha_i}A_i+r_i I_{\im E_i},\qquad
 r_i=\sum_{j>i}t^{\beta_j}.
\end{equation}
Here $r_m=0$, and otherwise $v(r_i)>\alpha_i$. The block on $\im P_m$ is zero.
Use any fixed real basis of each $\im E_i$. The normalized resolvent on that
block, at $s=ct^\alpha$, is
\[
 ct^\alpha\bigl(ct^\alpha I+t^{\alpha_i}A_i+r_iI\bigr)^{-1}.
\]
If $\alpha<\alpha_i$, factor out $ct^\alpha$; the residue is the identity.
If $\alpha>\alpha_i$, factor out $t^{\alpha_i}$; invertibility of the real
$A_i$ shows that the residue is zero. If $\alpha=\alpha_i$, the residue is
$c(cI+A_i|_{\im E_i})^{-1}$. These residue statements use only inversion of a
matrix with invertible real residue. Therefore the only crossovers are the
$\alpha_i$, and between them the shadows telescope to the prescribed $P_i$.

At $\alpha_j$ the later blocks together with $P_m$ sum to $P_j$. Under
$U_j,V_j$, this is the zero block $Q_j$ of $G_j$, while $\im E_j$ is its
complement. Consequently the full shadow is exactly
\eqref{eq:prescribed-shadow}. Since $p_{j-1}>p_j$, that function has different
endpoints and the crossover is nontrivial.

\emph{Step 4: impose irreducibility and then monomial rates.}
Apply Lemma~\ref{lem:completion} at a scale above $\alpha_m$. The resulting
$L$ has a positive rate on every ordered pair and still has finite support.
Each rate is its positive leading monomial times $1+\varepsilon_e$ with
$v(\varepsilon_e)>0$. Theorem~\ref{thm:leading} already proves that replacing
it by that leading monomial changes none of the shadows. The stronger relative
statement will be established in Theorem~\ref{thm:stability}.
\end{proof}

\begin{corollary}[A necessary and sufficient criterion]\label{cor:classification}
A finite collection of plateau projections and real crossover kernels is the
leading resolvent hierarchy of an irreducible finite surreal rate matrix if and
only if its plateaux form a strict nested stochastic-idempotent flag from $I$ to
rank one, and each crossover, on the preceding stochastic retract, is the
normalized resolvent of a real Markov generator whose final projection is the
following plateau. Any prescribed strictly ordered Hahn locations for the
crossovers can be used.
\end{corollary}

\begin{proof}
Necessity is Theorems~\ref{thm:flag} and~\ref{thm:effective}. Sufficiency is
Theorem~\ref{thm:general-realization}, transported into $\No$ by
\eqref{eq:surreal-embedding} whenever the chosen scales are surreal exponents.
\end{proof}

\subsection{The reversible boundary}

\begin{proposition}[Full-support reversible flags]\label{prop:reversible}
Suppose a surreal rate matrix is reversible with stationary probability $\pi$
and every entry of $\pi_0=\res\pi$ is strictly positive. Then every plateau is
the conditional-expectation matrix of a partition, weighted by $\pi_0$, and the
partitions are nested by coarsening. Conversely, every such finite nested
partition family is realizable by a reversible irreducible Hahn generator at
arbitrary prescribed scales.
\end{proposition}

\begin{proof}
Detailed balance gives
$\diag(\pi)R_L(s)=R_L(s)^T\diag(\pi)$. Taking residues shows that each plateau
$P$ is self-adjoint for the positive real weights $\pi_0$. It also preserves
$\pi_0$. The splitting proof shows that any transient state of a stochastic
idempotent has a zero column. Such a state cannot have positive mass under a
stationary probability, so there are no transient states. Each closed block is
$\one\nu_a$, and detailed balance forces
$\nu_a(j)=\pi_0(j)/\pi_0(C_a)$ on its class. This is precisely weighted
conditional expectation onto that partition. The range consists of observables
constant on each block; nested ranges give coarsening partitions.

Conversely use \eqref{eq:simple-inverse} for the weighted partition projections.
Every summand is self-adjoint for $\pi_0$, so detailed balance holds. The last
projection is $\one\pi_0$ with strictly positive entries. Expression
\eqref{eq:inverse-signs} then makes every off-diagonal entry strictly negative.
\end{proof}

The full-support condition cannot be removed. Example~\ref{ex:fuzzy} gives an
exactly reversible system whose limiting projection has a genuinely mixed
transient row. The broader stochastic-retract theorem is not just a reformulation
of the reversible partition case.

\section{Uniform relative stability at every scale}\label{sec:stability}

Resolvent perturbation estimates based on a norm of $(sI+L)^{-1}$ can lose
precision near a small spectral gap. Positivity allows a different statement:
control relative edge errors by relative \emph{entry} errors, simultaneously for
all positive $s$.

\begin{theorem}[All-parameter, entrywise relative stability]\label{thm:stability}
Let $L,L'$ have the same strongly connected allowed graph. Suppose
\[
 q'_e=q_e(1+\eta_e),\qquad v(\eta_e)\ge\delta>0
 \quad(e\in E).
\]
For every positive $s\in F$, and every $i,j$,
\begin{equation}\label{eq:relative-stability}
 v\!\left(\frac{R_{L'}(s)_{ij}}{R_L(s)_{ij}}-1\right)\ge\delta.
\end{equation}
More generally the same bound holds with $R_{L'}(s')$ in the numerator if
$s'=s(1+\eta_s)$ and $v(\eta_s)\ge\delta$. There is no dependence on a spectral
gap, on $s$, or on the valuation of the entry being compared.
\end{theorem}

\begin{proof}
For every forest $f$,
\[
 q'(f)(s')^{d-|f|}
 =q(f)s^{d-|f|}\prod_{e\in f}(1+\eta_e)(1+\eta_s)^{d-|f|}
 =q(f)s^{d-|f|}(1+\bigO{\delta}).
\]
Let $Z$ be either the denominator forest sum $D(s)$ or any marked numerator
$N(s)_{ij}$, and let $Z'$ be its perturbed counterpart. Every summand of $Z$ is
positive. If their minimum valuation is $\zeta$, then $v(Z)=\zeta$ without
cancellation, whereas $v(Z'-Z)\ge\zeta+\delta$. Hence
$Z'/Z=1+\bigO{\delta}$. Both numerator and denominator are nonzero, because the
graph is strongly connected. Dividing the two relative identities gives
\eqref{eq:relative-stability}; the inverse of $1+\bigO{\delta}$ is again
$1+\bigO{\delta}$.
\end{proof}

\begin{corollary}[Leading data suffice simultaneously]\label{cor:leading-equivalence}
Two positive rate systems with the same graph, edge valuations, and edge leading
coefficients have identical plateau projections, crossover kernels, effective
generators, and leading terms of all marked entries at every scale. The assertion
holds simultaneously over all positive resolvent parameters, not only at the
finitely many critical scales.
\end{corollary}

\begin{proof}
For finitely many edges, the relative discrepancies have a common positive
valuation lower bound. Apply Theorem~\ref{thm:stability}. The effective generators
are determined by the compressed kernels in \eqref{eq:effective-one}.
\end{proof}

\begin{example}[Sharpness of the valuation bound]\label{ex:sharp}
Take two states, $q_{12}=q_{21}=1$, and $s=1$. Then $R_L(1)_{12}=1/3$.
Change only $q_{12}$ to $1+\eta$, with $v(\eta)=\delta>0$. The new entry is
$(1+\eta)/(3+\eta)$, and its relative error is exactly
\[
 \frac{(1+\eta)/(3+\eta)}{1/3}-1=\frac{2\eta}{3+\eta}.
\]
Its valuation is $\delta$. Thus the exponent in
\eqref{eq:relative-stability} cannot be increased in general.
\end{example}

\begin{remark}[Why the hypotheses matter]
A new edge is not a relative perturbation of a zero edge. It may change the
hierarchy, unless an additional argument such as Lemma~\ref{lem:completion}
shows otherwise. Likewise, positive leading forest sums are essential: for a
signed sum, near cancellation can amplify relative errors. The theorem is not
an arbitrary-matrix resolvent condition-number bound.
\end{remark}

\section{A complete finite diagram has a rank-one realization}\label{sec:specialization}

Higher rank is genuine, but the finite nature of a forest calculation creates a
precise limitation on what it can detect. This section preserves not merely the
critical eigenvalue scales, but every marked entry's leading forest choices.

\subsection{Finite sign separation}

\begin{lemma}[Finite positive signs can be represented rationally]\label{lem:separation}
Let $W$ be a finite-dimensional rational vector space with an ordered-abelian-group
order, and let $x_1,\ldots,x_r$ be a finite family of strictly positive elements.
There exists a $\Q$-linear map $\ell:W\to\Q$ such that $\ell(x_i)>0$ for every
$i$. No assertion of positivity on all of $W_{>0}$ is made.
\end{lemma}

\begin{proof}
If the list is empty, take the zero functional. Otherwise choose a rational
basis of $W$ and regard the $x_i$ as rational vectors in
$\R^a$. Their convex hull cannot contain zero. Otherwise a convex representation
of zero can be chosen with a minimal number of nonzero coefficients. The
augmented vectors $(x_i,1)$ in that representation are linearly independent:
a dependence would allow the coefficients to be varied while preserving their
sum and the represented vector, until one coefficient became zero. Solving the
resulting independent rational linear system makes the minimal representation
rational. It would therefore express zero as a nonempty positive rational
combination of positive elements of the ordered group, a contradiction.

The convex hull is compact and excludes zero. If $y$ is its point of minimal
Euclidean norm, the usual nearest-point inequality gives
$y\cdot x_i\ge\|y\|^2>0$ for each $i$. A sufficiently close rational vector has
the same finitely many strict inequalities. Dot product with it defines the
required rational functional.
\end{proof}

This elementary finite separation principle is used as an input, not claimed as
new. Higher-rank Hahn specialization in more general tropical settings is studied
in~\cite{JS}.

\subsection{The marked-forest arrangement}

For every forest $f$ define an affine function on the ordered group,
\[
 \lambda_f(\alpha)=w(f)+(d-|f|)\alpha.
\]
All coefficient data belong to the finite-dimensional space
\[
 W=\Span_\Q\{\gamma_e:e\in E\}\subseteq\Gamma.
\]
If two forest functions have different slopes, their intersection is a point of
$W$. Let $T$ be the finite set of all such intersection points. We retain
intersections that are not active for the denominator, since some of them may
matter to a marked numerator. Forest functions with the same slope are ordered
by their constant terms, or coincide identically.

\begin{theorem}[Simultaneous rank-one realization of the leading diagram]
\label{thm:rank-one}
For any fixed finite strongly connected Hahn rate system there is a
$\Q$-linear map $\ell:W\to\Q$ such that the ordinary positive rates
\begin{equation}\label{eq:real-specialization}
 q_e(x)=a_e x^{\ell(\gamma_e)},\qquad x>0,
\end{equation}
reproduce its entire finite marked-forest leading diagram as $x\downarrow0$.
Specifically, the ordered cells and crossing points of the forest arrangement
correspond, and at corresponding locations they have the same minimizing forests
for the denominator and for every marked numerator. Consequently they have the
same rational amplitude formulas, plateau matrices, crossover matrices, and
effective generators.

At a crossing $\alpha\in T$, its numerical location is $\ell(\alpha)$.
On an open cell, the corresponding affine exponent formulas are obtained by
applying $\ell$ to their constant terms and replacing the scale variable by a
real variable in the corresponding cell. This is not a global order embedding
of $\Gamma$. After a change $x=y^M$, all exponents in the rates can be chosen
integral, possibly negative.
\end{theorem}

\begin{proof}
Form a finite list of positive elements of $W$ consisting of all positive
differences between distinct points of $T$, and all positive differences
$w(f)-w(g)$ between equal-slope forest functions. Optional finitely many positive
edge-valuation differences or other prescribed comparisons may be added.
Apply Lemma~\ref{lem:separation}. It preserves the strict order of the crossing
points and the signs of equal-slope differences. Equalities are automatically
preserved by linearity.

For two different-slope functions, their difference is an integer multiple of
$\alpha-\theta$, where $\theta\in T$ is their intersection. Thus its sign on
any cell, or at a crossing, is determined by the ordering relative to $\theta$.
That ordering is preserved. Together with the equal-slope comparisons, this
shows that the total comparison pattern of all forest functions is preserved
in each cell and at each crossing. Restricting the family to any set
$r_f(i)=j$ preserves its minimizers as well.

For the rates \eqref{eq:real-specialization}, a forest term at the numerical
observation scale $s=cx^a$ has exponent
$\ell(w(f))+(d-|f|)a$ and coefficient $A(f)c^{d-|f|}$. Finite sums of positive
monomials are asymptotic to the sum of the terms with smallest exponent. Since
the minimizing forests and their coefficients agree, all amplitude formulas in
Theorem~\ref{thm:leading} agree. In particular their matrix residues, plateau
splittings, and the reconstructed real generators agree.

There are finitely many rational edge exponents. A common positive denominator
$M$ makes the exponents of $q_e(y^M)$ integral. If desired, multiplying all rates
by the same power of $y$ makes them nonnegative, at the cost of translating every
observation-scale exponent by that same amount.
\end{proof}

\begin{remark}[Precisely what is and is not preserved]
An arbitrary $\alpha\in\Gamma$ need not belong to $W$, so an expression
$\ell(\alpha)$ is not used for every possible scale. The finite arrangement
provides the correspondence: select a real scale in the matching open cell, or
use $\ell(\alpha)$ at a crossing. The marked exponent
$m_{ij,\alpha}-m_\alpha$ is a piecewise affine expression. Its combinatorial
formula and amplitude are preserved; its value as an element of $\Gamma$ is
not literally identified with a real number unless the relevant arguments are
in the domain of $\ell$. The theorem preserves leading data, not complete Hahn
series, all possible inequalities, or the entire ordered group.
\end{remark}

\begin{example}[Why the finite qualification is sharp]\label{ex:rank-obstruction}
In $\Gamma=\Q^2$ with lexicographic order put
$\delta=(0,1)$ and $\beta=(1,0)$. Then
\[
 0<N\delta<\beta\qquad\text{for every positive integer }N.
\]
No real linear functional can preserve all these inequalities while sending
$\delta$ to a positive number: it would require the finite real number
$\ell(\beta)$ to exceed $N\ell(\delta)$ for every $N$. A finite family of
inequalities can be preserved, for example by sending $\delta$ to $1$ and
$\beta$ to a sufficiently large integer. Theorem~\ref{thm:rank-one} uses exactly
such a finite list, not its infinite continuation.
\end{example}

\section{The surcomplex spectrum and the forest convexity law}\label{sec:spectrum}

Because $K=\C((t^\Gamma))$ is algebraically closed, the eigenvalues of $L$ lie
in a surcomplex Hahn workspace. The tree coefficient $\sigma_d>0$ in
\eqref{eq:forest-det} shows that zero has algebraic multiplicity one. Denote the
other eigenvalues, with algebraic multiplicity, by
$\lambda_1,\ldots,\lambda_d$, and put $\zeta_r=v(\lambda_r)$.

\subsection{Initial polynomials without a real Newton polygon}

\begin{lemma}[Valuations and initial eigenvalue factors]\label{lem:initial-spectrum}
For every $\alpha\in\Gamma$,
\begin{equation}\label{eq:sum-min}
 m_\alpha=\sum_{r=1}^d\min(\alpha,\zeta_r).
\end{equation}
If $\rho_r=\res(t^{-\alpha}\lambda_r)$ for the indices with
$\zeta_r=\alpha$, then
\begin{equation}\label{eq:initial-factorization}
 D_\alpha(c)=C_\alpha\,
 c^{\#\{r:\zeta_r>\alpha\}}
 \prod_{\zeta_r=\alpha}(c+\rho_r),
\end{equation}
where $C_\alpha\ne0$ is the product of the leading coefficients of the faster
factors. Consequently the critical scales are exactly the distinct nonzero
eigenvalue valuations.
\end{lemma}

\begin{proof}
As a polynomial over $K$,
$D(s)=\prod_r(s+\lambda_r)$. Replace $s$ by $t^\alpha z$, where $z$ is an
indeterminate. The minimum coefficient valuation of a product of nonzero
polynomials is the sum of those minima: after normalizing each factor to have
minimum zero, the product of their nonzero residue polynomials remains nonzero
in the integral domain $\C[z]$.

The normalized residue of $t^\alpha z+\lambda_r$ is respectively a nonzero
constant, $z+\rho_r$, or $z$, according as $\zeta_r<\alpha$,
$\zeta_r=\alpha$, or $\zeta_r>\alpha$. This proves both formulas. If no
$\zeta_r$ equals $\alpha$, the initial polynomial is a monomial. If at least
one does, its largest and smallest nonzero powers differ, since all those
$\rho_r$ are nonzero. Hence it is not a monomial. By
\eqref{eq:initial-polynomials}, being nonmonomial is exactly criticality.
\end{proof}

\subsection{A characteristic polynomial for each effective chain}

\begin{theorem}[Effective characteristic and amplitude identity]\label{thm:char}
At a critical scale $\alpha$, let $k_-,k_+,p,q$ be as in
Theorem~\ref{thm:flag}. The effective generator satisfies the exact real polynomial
identity
\begin{equation}\label{eq:effective-char}
 \det(cI_p+G_\alpha)=\frac{cD_\alpha(c)}{b_{k_-}}.
\end{equation}
Exactly $p-q$ nonzero eigenvalues of $L$ have valuation $\alpha$, counting
algebraic multiplicity. Their normalized leading coefficients are exactly the
nonzero eigenvalues of $G_\alpha$, again counting multiplicity. In particular
these coefficients have strictly positive real parts, and every nonzero
surcomplex eigenvalue satisfies
\begin{equation}\label{eq:realpart-scale}
 v(\Rea\lambda)=v(\lambda).
\end{equation}
\end{theorem}

\begin{proof}
The trace relation in \eqref{eq:profile-properties} gives
\begin{equation}\label{eq:trace-log}
 \tr K_\alpha(c)
 =\frac{\sum_{k\in I_\alpha}(n-k)b_kc^{d-k}}{D_\alpha(c)}
 =1+c\frac{D_\alpha'(c)}{D_\alpha(c)}.
\end{equation}
On the other hand, cyclicity of trace and $VU=I$ give
\[
 \tr K_\alpha(c)=\tr\bigl(c(cI_p+G_\alpha)^{-1}\bigr)
 =c\frac{d}{dc}\log\det(cI_p+G_\alpha).
\]
Therefore the logarithmic derivatives of
$\det(cI_p+G_\alpha)$ and $cD_\alpha(c)$ agree for positive $c$. Their ratio is
a constant rational function. The former is monic of degree $p$; the latter
has degree $1+d-k_-=p$ and leading coefficient $b_{k_-}$. This proves
\eqref{eq:effective-char} as a polynomial identity.

The smallest power of $D_\alpha$ is $d-k_+=q-1$, and the largest is $p-1$.
Comparing with \eqref{eq:initial-factorization}, there are $q-1$ slower
nonzero eigenvalues and $p-q$ at the current scale. Cancelling the $q$ zero
factors in \eqref{eq:effective-char} identifies the other eigenvalues with the
$\rho_r$ of that factorization. Nonzero eigenvalues of a real row Laplacian have
strictly positive real parts, as proved in Corollary~\ref{cor:effective-semigroup}.
Thus the leading real part of every normalized $\lambda_r$ is nonzero and
positive, proving \eqref{eq:realpart-scale}.
\end{proof}

There is no assertion here that a nonreversible spectrum is real. The leading
amplitudes can occur in nonreal conjugate pairs, as in
Example~\ref{ex:prescribed}. Nor is diagonalizability at nonzero eigenvalues
needed for the determinant argument.

\begin{theorem}[Discrete convexity and all eigenvalue valuations]
\label{thm:convexity}
The forest profile satisfies
\begin{equation}\label{eq:convexity}
 h_1-h_0\ \le\ h_2-h_1\ \le\ \cdots\ \le\ h_d-h_{d-1}.
\end{equation}
These $d$ increments are precisely the nonzero eigenvalue valuations of $L$ in
nondecreasing order, including repetitions. At any critical scale,
$I_\alpha$ consists of every integer from $k_-$ through $k_+$; no intermediate
forest size is missing from the initial polynomial.
\end{theorem}

\begin{proof}
Apply the classical forest determinant formula to the real effective row
Laplacian $G_\alpha$. Its zero eigenvalue has multiplicity $q$, so the largest
number of edges in one of its in-forests is $p-q$. Every smaller edge count
occurs by deleting edges of such a forest. All associated coefficients of
$\det(cI_p+G_\alpha)$ are positive. Thus this polynomial has strictly positive
coefficients at all degrees from $q$ to $p$, inclusive. Equation
\eqref{eq:effective-char} implies that $D_\alpha$ has a positive coefficient
at every degree from $q-1$ to $p-1$. Its definition forces every integer from
$k_-$ to $k_+$ to belong to $I_\alpha$.

As the scale increases from the initial to the final plateau, the active indices
run from $0$ to $d$. Any jump from $k_-$ to $k_+$ therefore has all intermediate
points on the same supporting line, and
\[
 h_k-h_{k-1}=\alpha\qquad(k_-<k\le k_+).
\]
Successive jumps occur at strictly increasing scales. This proves
\eqref{eq:convexity}. The multiplicity statement in Theorem~\ref{thm:char}
then identifies the increments with the eigenvalue valuation list.
\end{proof}

\begin{remark}[Relation to optimal-forest spectral formulas]
In a rank-one small-noise model, forest costs are sums of edge exponents, and
successive optimal costs give familiar eigenvalue exponents in the
Freidlin--Wentzell tradition; see the discussion in~\cite{GC}. The present proof
works in an arbitrary ordered group and simultaneously identifies the full
real effective characteristic polynomial. It does not claim that the underlying
rank-one optimal-forest idea is new.
\end{remark}

\section{Worked examples}\label{sec:examples}

\subsection{A nonreversible three-state system}

Let $\delta>0$, write $\tau=t^\delta$, and consider
\begin{equation}\label{eq:example-three}
 L=\begin{pmatrix}
 2&-2&0\\
 -1&1+3\tau&-3\tau\\
 -\tau&0&\tau
 \end{pmatrix}.
\end{equation}
The graph is strongly connected, but is not reversible: there is an edge
$2\to3$ without its reverse. Exact expansion gives
\[
 D(s)=s^2+(3+4\tau)s+9\tau+3\tau^2.
\]
The forest profile is
\[
 (h_0,h_1,h_2)=(0,0,\delta),\qquad (b_0,b_1,b_2)=(1,3,9).
\]
Thus the nonzero eigenvalue scales are $0$ and $\delta$. Between them and after
them, respectively, the plateau projections are
\begin{equation}\label{eq:three-projections}
 P=\begin{pmatrix}
 1/3&2/3&0\\1/3&2/3&0\\0&0&1
 \end{pmatrix},\qquad
 Q=\one\begin{pmatrix}1/9&2/9&2/3\end{pmatrix}.
\end{equation}
At the slower scale,
\begin{equation}\label{eq:three-crossover}
 K_\delta(c)=\frac{cP+3Q}{c+3}.
\end{equation}
Use
\[
 U=\begin{pmatrix}1&0\\1&0\\0&1\end{pmatrix},\qquad
 V=\begin{pmatrix}1/3&2/3&0\\0&0&1\end{pmatrix}.
\]
Compression reconstructs
\[
 G_\delta=\begin{pmatrix}2&-2\\-1&1\end{pmatrix}.
\]
Its nonzero eigenvalue is $3$, so the slow eigenvalue of $L$ has leading term
$3t^\delta$. The fast eigenvalue has leading term $3$. The characteristic identity
reads $\det(cI_2+G_\delta)=c(c+3)=cD_\delta(c)/3$, because
$D_\delta(c)=3c+9$.

One rare entry illustrates the marked formula. For $0<\alpha<\delta$, the
plateau entry $P_{13}$ is zero. The leading numerator term for $R_{13}$ is
$6\tau$, whereas the denominator has leading term $3ct^\alpha$. Therefore
\[
 R_L(ct^\alpha)_{13}
 =\frac{2}{c}\,t^{\delta-\alpha}
     \bigl(1+\text{an element of positive valuation}\bigr).
\]
The vanishing shadow does not mean the transition is absent; it has a definite
smaller leading scale and amplitude.

\subsection{A rank-two hierarchy with a tied fast scale}\label{subsec:rank-two}

In lexicographically ordered $\Q^2$, let $\delta=(0,1)$, $\beta=(1,0)$,
$\tau=t^\delta$, and $\zeta=t^\beta$. Then $\zeta\ll\tau^N$ for every positive
integer $N$. Let $P$ average each pair $\{1,2\}$ and $\{3,4\}$, and let
$J=\one\one^T/4$. The inverse construction gives
\[
 H=\tau(I-P)+\zeta(P-J).
\]
The rate within a pair is $\tau/2-\zeta/4>0$, and each cross-pair rate is
$\zeta/4$. This is a reversible complete graph. Its nonzero eigenvalues are
$\tau,\tau,\zeta$, and
\[
 D(s)=(s+\tau)^2(s+\zeta).
\]
Consequently
\[
 (h_0,h_1,h_2,h_3)=(0,\delta,2\delta,2\delta+\beta),
 \quad(b_0,b_1,b_2,b_3)=(1,2,1,1).
\]
The two crossovers are
\[
 K_\delta(c)=P+\frac{c}{c+1}(I-P),\qquad
 K_\beta(c)=J+\frac{c}{c+1}(P-J).
\]
The fast transition loses two dimensions at once. Tied scales cause no
ambiguity and require no arbitrary tie breaking. A numerical monomial
specialization can reproduce this entire finite diagram, but cannot preserve
$\zeta\ll\tau^N$ for every $N$; this is the distinction in
Theorem~\ref{thm:rank-one}.

\subsection{Reversibility does not eliminate transient mixtures}

\begin{example}[A fuzzy plateau in a reversible system]\label{ex:fuzzy}
Let $p,q>0$ be real with $p+q=1$, and let $\tau=t^\delta$, $\delta>0$. Put
\[
 L=\begin{pmatrix}
 \tau&0&-\tau\\
 0&\tau&-\tau\\
 -p&-q&1
 \end{pmatrix}.
\]
It is reversible for the exact stationary law
\[
 \pi=\frac{1}{1+\tau}(p,q,\tau),
\]
since the two edgewise detailed-balance identities are immediate. Its determinant
factor is
$D(s)=(s+\tau)(s+1+\tau)$. At scales $0<\alpha<\delta$, its plateau is
\[
 P=\begin{pmatrix}1&0&0\\0&1&0\\p&q&0\end{pmatrix}.
\]
The third row is a genuine probability mixture between the two surviving states.
The canonical retract uses the first two coordinate rows for $V$ and the first
two columns of $P$ for $U$. Its slow effective generator is
\[
 G_\delta=\begin{pmatrix}q&-q\\-p&p\end{pmatrix}.
\]
The residue stationary law $(p,q,0)$ is not strictly positive. Thus this example
satisfies exact reversibility while lying outside the full-support conclusion
of Proposition~\ref{prop:reversible}.
\end{example}

\subsection{Prescribed nonreal dynamics and an invisible positivity correction}

\begin{example}[Why the general inverse theorem needs its correction]
\label{ex:prescribed}
For clarity take a rank-one parameter $t$ with $v(t)>0$, measuring exponents in
units of $v(t)$. Prescribe a fast generator that cycles through the first three
states and leaves the fourth fixed:
\[
 A_1=\begin{pmatrix}
 1&-1&0&0\\0&1&-1&0\\-1&0&1&0\\0&0&0&0
 \end{pmatrix}.
\]
Its zero-eigenvalue projection is
\[
 P=UV,\quad
 U=\begin{pmatrix}1&0\\1&0\\1&0\\0&1\end{pmatrix},\quad
 V=\begin{pmatrix}1/3&1/3&1/3&0\\0&0&0&1\end{pmatrix}.
\]
Prescribe the slow effective generator
\[
 G_2=\begin{pmatrix}2&-2\\-1&1\end{pmatrix},\qquad A_2=UG_2V,
\]
and place it at scale $2$. The naive sum $A_1+t^2A_2$ is not a row Laplacian:
its $(1,3)$ entry is $2t^2/3>0$. The intermediate correction gives
\begin{equation}\label{eq:corrected-four}
 H=A_1+t(I-P)+t^2A_2.
\end{equation}
Now the offending entry is $-t/3+2t^2/3<0$, and all the other off-diagonal signs
are correct by the proof of Theorem~\ref{thm:general-realization}.

The exact nonzero characteristic factor is
\begin{equation}\label{eq:corrected-four-char}
 D(s)=(s+3t^2)\bigl((s+t)^2+3(s+t)+3\bigr).
\end{equation}
Thus the eigenvalues are
\[
 0,\quad 3t^2,\quad
 t+\frac32+\frac{\sqrt3}{2}i,\quad
 t+\frac32-\frac{\sqrt3}{2}i.
\]
This supplies a genuine surcomplex pair at the fast scale. The complete
resolvent can be written
\begin{equation}\label{eq:corrected-four-resolvent}
 R_H(s)=s((s+t)I+A_1)^{-1}(I-P)
       +U\,s(sI_2+t^2G_2)^{-1}V.
\end{equation}
At scale $0$ its shadow is $c(cI_4+A_1)^{-1}$; at scale $2$ it is
$U c(cI_2+G_2)^{-1}V$. At the inserted scale $1$, for every real $c>0$, the
shadow is just $P$. A new leading \emph{edge} scale has repaired positivity
without becoming a new leading \emph{resolvent} scale.

The final projection is
$\one(1/9,1/9,1/9,2/3)$, and this particular corrected matrix is already
irreducible. Multiplication of \eqref{eq:corrected-four-resolvent} by $sI+H$
provides a direct exact check independent of the asymptotic argument.
\end{example}

\section{Computation, verification, and limitations}\label{sec:computation}

\subsection{A finite exact procedure}

The theory gives a direct algorithm once the leading edge data are available.
Enumerate in-forests by assigning to each vertex either no outgoing edge or one
allowed successor, and reject assignments containing a directed cycle. For each
surviving forest, compute its root map, valuation cost, and positive amplitude.
Grouping by edge count gives $h_k,b_k,B_k$; grouping additionally by the mark
$r_f(i)=j$ gives the rare-entry data.

The increments $h_k-h_{k-1}$ already give the critical scales with multiplicity,
by Theorem~\ref{thm:convexity}. At a critical scale form
$N_\alpha(c)/D_\alpha(c)$ and the preceding projection $P$. Find its closed
classes to construct $U,V$, then use
\[
 G_\alpha=(VK_\alpha(1)U)^{-1}-I.
\]
Thus a single real matrix inverse per scale reconstructs the effective generator
once the positive forest coefficients are known. Computing surcomplex eigenvectors
is unnecessary.

This enumeration is finite, not an efficiency claim. There are at most $n^n$
parent assignments in a complete directed graph, and the number of forests grows
rapidly. The exact comparison contract is also significant: one must be able to
add and compare the input valuations, perform rational division in their span,
and manipulate the real leading coefficients exactly. Arbitrary real coefficient
names need not supply decidable equality; arbitrary surreal names need not expose
the leading exponent. A mathematical finiteness theorem is not a uniform
algorithm on every possible representation of a surreal number.

\subsection{Reproducible checks included in the archive}

The accompanying \texttt{code/verify.py} uses Python's exact rational numbers and
SymPy, with no floating-point tolerances and no external data. The recorded run
used seed \texttt{20260922}. It tested 48 complete directed graphs on two through
five vertices, with positive rational leading coefficients and valuations in
lexicographically ordered $\Q^2$. Across these instances it enumerated 17,280
forests and checked 105 critical kernels.

The checks include convexity of the optimal forest profile; equality of its
increments with the critical scale multiset; stochasticity, idempotence and
nesting of all plateaux; the canonical stochastic retract identities; the signs
and ranks of the reconstructed effective generators; and equality with their
resolvents at the exact rational parameters $c=1,2,3$. These checks use forests
with tied costs as well as distinct costs.

Separate symbolic checks verify the three-state determinant and crossover,
the exact projection-flag inverse, the rank-one irreducible-completion identity,
a four-state matrix-forest identity, the sharp relative-error example, and the
prescribed-crossover construction in Example~\ref{ex:prescribed}, including the
absence of an extra crossover at its positivity-correction scale. The resulting
\texttt{data/verification.json} records \texttt{PASS} and the individual instance
sizes. These checks do not establish the general theorems: in particular they do
not replace the arbitrary-rank separation proof or the positivity argument for
all compatible effective generators. No Lean verification is claimed.

\subsection{What has actually been settled here}

The forward and inverse theorems settle an internally stated classification
problem: the necessary and sufficient conditions for a \emph{finite leading
resolvent hierarchy} to arise from positive surreal rates. They include the
realizability of arbitrary compatible reduced dynamics, rather than merely an
ordered eigenvalue list. This is a definite theorem with a constructive proof,
not a suggestion to investigate the problem later.

The marked formula and the relative bound give two different levels of
information. The unnormalized real shadow describes effective stochastic motion.
The marked entry's valuation and amplitude describe transitions too rare to
appear in that shadow. Relative stability controls both uniformly over the
observation parameter. The characteristic identity, in turn, explains why the
same forest data describe the leading surcomplex relaxation amplitudes.

The specialization theorem gives a complementary negative conclusion:
this finite information does not distinguish higher rank from a suitably chosen
rank-one model. Higher rank remains meaningful for an unbounded family of scale
comparisons, for complete series rather than leading terms, and for data whose
size is not fixed in advance. The finite theorem must not be extrapolated to
those settings.

\subsection{Remaining directions, not results claimed in this paper}

The normalized resolvent does not determine a theory of surreal-valued sample
paths, stochastic integration, or a globally defined surreal heat semigroup.
Constructing such objects requires additional choices of convergence and
summability structures. Likewise, no result here covers countably infinite state
spaces, time-inhomogeneous rates, uniform estimates as $n\to\infty$, or a full
hierarchy of large-deviation transition events.

There are natural further questions. One is to develop efficient exact algorithms
for the marked-forest diagram without exhaustive enumeration. Another is to
identify conditions under which the real effective semigroups arise from a
specified, independently defined Hahn-summable evolution. A third is to ask what
additional finite data, beyond leading resolvent terms, must be retained to
obstruct some rank-one specializations. These questions are separated from the
proved realization theorem and are not presented as solved.

\section{Conclusion}

Positive surreal transition rates admit a particularly rigid finite theory once
the observable is fixed. Their all-scale resolvent shadows are exactly finite
compatible cascades of real Markov dynamics on nested stochastic retracts. The
inverse construction shows that no unmentioned reversibility or partition
condition is hidden in this description. Its positivity corrections can occupy
new edge scales without creating new leading dynamical scales.

The same finite forest algebra supplies uniform relative perturbation control,
rare-entry asymptotics, the effective characteristic polynomial, and an exact
rank-one realization of the whole finite leading diagram. This is simultaneously
a constructive theory of multiscale surreal Markov matrices and a precise limit
on what finite leading-order observations can reveal about the rank of the
surreal scale group.

\appendix
\section{Proof dependencies and conventions at a glance}

The only imported structural results about generalized numbers are Conway normal
form, the standard Hahn-field arithmetic, and real/algebraic closedness for the
stated divisible groups. The only imported combinatorial matrix identity is the
directed matrix-forest theorem. The stochastic-idempotent structure theorem and
the finite sign-separation lemma are reproved. All realization and stability
arguments are then finite algebra.

\begin{center}
\begin{tabular}{p{0.33\textwidth}p{0.59\textwidth}}
\toprule
Conclusion & Main dependencies\\
\midrule
Leading and marked shadows & Positive finite forest sums; Proposition~\ref{prop:forest}.\\
Nested stochastic flag & Resolvent identity; trace of a forest matrix.\\
Effective real generator & Stochastic retract; compressed resolvent identity.\\
Prescribed-crossover realization & Block support of compatible generators; sign-cover implication \eqref{eq:sign-cover}; intermediate corrections.\\
Irreducible completion & Exact rank-one inverse identity; rank-one final projection.\\
Relative stability & No cancellation in positive forest sums.\\
Rank-one specialization & Finite affine forest arrangement; finite rational sign separation.\\
Surcomplex amplitudes & Polynomial factorization over $K$; effective characteristic identity.\\
Forest convexity & Positivity of every attainable forest coefficient of the real effective generator.\\
\bottomrule
\end{tabular}
\end{center}

There is no circular reliance between the inverse and stability theorems. The
inverse theorem's monomial replacement assertion already follows from the
leading-forest formula. The later stability theorem strengthens it but is not
needed to construct the hierarchy. Similarly, convexity is proved after the
effective generator exists; convexity is convenient for computation but was not
assumed in constructing that generator.

The time parameter is always distinguished from the scale exponent:
$s=ct^\alpha$ is a Hahn killing rate, $c$ is an ordinary positive real amplitude,
and $u$ in the effective semigroup is an ordinary nonnegative real time. The
valuation is additive, so larger $\alpha$ means smaller $s$. Eigenvalues refer
to the row Laplacian $L$, not to the conventional generator $-L$; nonzero real
parts therefore have the positive sign.

\begin{thebibliography}{99}

\bibitem{Conway}
J.~H. Conway,
\emph{On Numbers and Games}, second edition,
A K Peters, 2001.
Standard background for surreal numbers and normal forms.

\bibitem{Gonshor}
H. Gonshor,
\emph{An Introduction to the Theory of Surreal Numbers},
London Mathematical Society Lecture Note Series 110,
Cambridge University Press, 1986.

\bibitem{CA}
P. Chebotarev and R. Agaev,
Forest matrices around the Laplacian matrix,
\emph{Linear Algebra and its Applications} \textbf{356} (2002), 253--274.
\doi{10.1016/S0024-3795(02)00388-9}.
Author-deposited version:
\url{https://arxiv.org/abs/math/0508178}.

\bibitem{GC}
T. Gan and M. Cameron,
A graph-algorithmic approach for the study of metastability in Markov chains,
\emph{Journal of Nonlinear Science} \textbf{27} (2017), 927--972.
\doi{10.1007/s00332-016-9355-0}.
\url{https://arxiv.org/abs/1607.00078}.

\bibitem{YZ}
G. Yin and H. Zhang,
Singularly perturbed Markov chains: Limit results and applications,
\emph{Annals of Applied Probability} \textbf{17} (2007), no.~1, 207--229.
\doi{10.1214/105051606000000682}.
\url{https://arxiv.org/abs/math/0703017}.

\bibitem{DGU}
M. Dyer, C. Greenhill and M. Ullrich,
\emph{Structure and eigenvalues of heat-bath Markov chains},
author-hosted manuscript dated 10 April 2014, especially Section~4.
\url{https://web.maths.unsw.edu.au/~csg/papers/DGU-heatbath.pdf}.

\bibitem{BK}
A. Bendikov and P. Krupski,
On the spectrum of the hierarchical Laplacian,
\emph{Potential Analysis} \textbf{41} (2014), 1247--1266.
\doi{10.1007/s11118-014-9409-6}.
\url{https://arxiv.org/abs/1308.4883}.

\bibitem{NV}
R. Nabben and R.~S. Varga,
Generalized ultrametric matrices---a class of inverse $M$-matrices,
\emph{Linear Algebra and its Applications} \textbf{220} (1995), 365--390.
\doi{10.1016/0024-3795(94)00086-S}.
Author-hosted copy:
\url{https://www.math.kent.edu/~varga/pub/paper_212.pdf}.

\bibitem{JS}
M. Joswig and B. Smith,
Convergent Hahn series and tropical geometry of higher rank,
\emph{Journal of the London Mathematical Society}
\textbf{107} (2023), no.~4, 1450--1481.
\doi{10.1112/jlms.12716}.
\url{https://arxiv.org/abs/1809.01457}.

\bibitem{Repo}
V. Reshetnikov,
\emph{Surreal}, repository snapshot
\texttt{4cf691c7d951e037739d32d9f5c387dcce724f3c},
examined 22 September 2026.
\url{https://github.com/VladimirReshetnikov/Surreal}.
Targeted comparison: \texttt{docs/README.md},
\texttt{docs/surcomplex/spectral-theory/README.md},
the \texttt{docs/new} inventory, and indexed source search.
See the audit accompanying this manuscript for the unexpanded-archive limitation.

\end{thebibliography}
\end{document}
